% from main.tex
%%%%%%%%%%  scan idx 100  |  folio 87  %%%%%%%%%%
Nous dirons que \(\mathcal C_0\) est \emph{localement quasi-relevable dans} \(\mathcal C\) si la condition plus faible suivante est réalisée : pour tout \(X\in\Ob S\), et tout épimorphisme \(u:E\longrightarrow F\) de \(\mathcal C_X\) avec \(F\in\Ob\mathcal C_{0,X}\), l'ensemble des \(f:U\to X\) de \(S_{/X}\) tels qu'il existe \(v:F'\longrightarrow f^*E\) avec \(F'\in\Ob\mathcal C_{0,U}\) tel que \(f^*(u)v\) soit un épimorphisme, est un raffinement de \(X\). Cette condition équivaut à la suivante : pour tout diagramme de \(\mathcal C_X\)
\[
\begin{tikzcd}
& H \arrow[d]\\
E \arrow[r,"u"] & F ,
\end{tikzcd}
\]
où \(u\) est un épimorphisme, et \(H\in\Ob\mathcal C_0\), l'ensemble des \(U\) au-dessus de \(X\) tels qu'il existe un diagramme commutatif
\[
\begin{tikzcd}
G \arrow[r,"v"] \arrow[d] & H_U \arrow[d]\\
E_U \arrow[r,"u_U"] & F_U
\end{tikzcd}
\]
avec \(G\in\Ob\mathcal C_0\) et \(v\) un épimorphisme, est un raffinement de \(X\).

Si \(S\) est un site chaotique, par exemple le site ponctuel (un objet, un morphisme), nous omettrons l'adverbe «~localement~» dans les définitions précédentes.

\subsection*{1.3. Exemples}

\textbf{a)} Soient \(S\) un site ponctuel, \(A\) un anneau, \(\mathcal C\) la catégorie des \(A\)-modules à gauche. La sous-catégorie pleine \(\mathcal C_0\) de \(\mathcal C\) formée des \(A\)-modules (à gauche) projectifs de type fini est stable par noyau d'épimorphisme et relevable dans \(\mathcal C\)~; un \(A\)-module est de \(\mathcal C_0\)-type fini si et seulement s'il est de type fini au
